\par\medskip
\begin{center}\textls[100]{c) Theorie der Modulfunktion.}\end{center}
\noindent\hspace*{\parindent}Wir müssen in diesem Abschnitt den streng historischen Standpunkt, den
innezuhalten wir bisher bemüht waren, etwas lockern und uns auf eine Exposition
des Inhaltes der hier in Betracht kommenden, zum Teil \textls[100]{sicher} aus
sehr viel späterer Zeit stammenden, Aufzeichnungen beschränken. Aber auch
diese Aufgabe ist insofern erschwert, als die gedachten Aufzeichnungen nicht
erkennen lassen, von welcher Seite her \textsc{Gauss} den Zugang zu den darin enthaltenen
Resultaten genommen hat. Nur soviel scheint festzustehen, dass er den in
der Tagebuchnotiz 109. genannten »gegenseitigen Zusammenhang zwischen den
unendlich vielen Bestimmungen des agM.« nicht auf elementarem Weg, d. h.
nicht von dem Algorithmus des agM. ausgehend gefunden\footnote{Ein solcher \textls[100]{elementarer} Weg ist übrigens bis heute noch nicht entdeckt.}, sondern dass er
sich dabei der Eigenschaften der Funktionen \(p,q,r\) bedient hat.

Die in Betracht kommenden Aufzeichnungen sind die folgenden:

1) Das aus der Scheda Ae, Juli 1800, stammende Stück VI, Werke X, 1,
S. 207, 208.

2) Die Notiz aus Scheda Af, 1801, Abschnitt VII, Werke X, 1, S. 209
bis 212.

3) Der Abschnitt IX, Werke X, 1, S. 217—231, insbesondere der auf
den Seiten 218—226 wiedergegebene Inhalt zweier, wohl um 1825 geschriebener
Zettel.

4) Die Stelle Werke III, S. 386, Zeile 5—1 v. u., die \textsc{Gauss} auf der
letzten Einbandseite seines Handexemplars der \textit{Disquisitiones arithmeticae} eingetragen hat.

5) Die Notiz [3] Werke VIII, S. 99—102, die in dem Handbuch 25, Bf,
S. 25 aufgezeichnet ist und wohl aus dem Jahre 1815 stammen dürfte.

6) Eine Aufzeichnung in der Scheda An (teilweise veröffentlicht Werke
VIII, S. 103), die wir weiter unten im Texte wiedergeben, und die wohl auf
1805 anzusetzen ist.\par
\clearpage\SourceLeaf{179}\ArithmeticHead{94}{L}{L. SCHLESINGER, ÜBER GAUSS’ FUNKTIONENTHEORETISCHE ARBEITEN.}
\indent 7) Die Werke III, S. 470—480 abgedruckte Abhandlung aus Handbuch
16, Bb, S. 137—145, die \textsc{Gauss} aus dem August 1827 datiert; namentlich der
Art. [12], Werke III, S. 477, 478 (S. 142, 143 des Handbuchs)\footnote{Wir zitieren diese Stücke im Folgenden mit der ihnen in der obigen Übersicht beigelegten Nummer.}.

Mit der Tagebuchnotiz 109. im unmittelbaren Zusammenhange steht die
Aufzeichnung Art. 3 von 3), wo es heisst (Werke X, 1, S. 219): »Die agMittel
gestalten sich anders, wenn man für ein \(b,b_1,b_2,\ldots\) den negativen Wert wählt;
doch sind alle Resultate in folgender Form begriffen:
\begingroup\makeatletter\tagsleft@true\let\veqno\leqno\makeatother
\begin{equation*}
\frac1{(\mu)}=\frac1\mu+\frac{4ik}\lambda\text{\raisebox{.3ex}{.}«}\tag{10}
\end{equation*}
\endgroup
\noindent Der Wortlaut dieser Aussage, wie auch das derselben hinzugefügte »Beispiel
für einen imaginären Wert des agM.« legt es nahe anzunehmen, dass \textsc{Gauss}
den Fall vor Augen hatte, wo \(\mu\) bezw. \(\lambda\) die positiven Werte der agM. für
die positiven Zahlen \(a>b\), bezw. \(a\) und \(c=(a^2-b^2)^{\frac12}\) bedeuten, während \((\mu)\)
so zustande kommt, dass für eines der geometrischen Mittel der negative Wert
genommen wird. Jedenfalls stellt, wie wir sehen werden, das durch (10) gegebene
\((\mu)\) nicht den allgemeinsten Wert des agM. zwischen \(a,b\) dar. Über
die Art und Weise, wie \textsc{Gauss} die Formel (10) hergeleitet haben mag, sind
wir auf Vermutungen angewiesen. Es liegt ja nahe, dass die Bedeutung von
\(\frac1\mu\) und \(\frac i\lambda\) als Perioden der elliptischen Funktion \(\operatorname{Sl}\varphi\) ihn dazu geführt hat,
anzunehmen, dass \(\frac1{(\mu)}\) linear homogen mit ganzzahligen Koeffizienten durch
\(\frac1\mu\) und \(\frac i\lambda\) darstellbar sein müsste. Der handschriftliche Nachlass bietet allerdings
keinen Anhaltspunkt hierfür. \textsc{Schering} deutet (Werke III, S. 377, 378)
einen Weg an, wie man mit Hilfe von Kontinuitätsbetrachtungen zu dieser
Formel gelangen könnte. Nach dem Inhalt der beiden Zettel 3) ist es aber
wahrscheinlicher, dass \textsc{Gauss} die in Rede stehende Formel als besondern Fall
des allgemeinen, in der Tagebuchnotiz 109. angezeigten Zusammenhangs zwischen
den unendlich vielen Bestimmungen des agM., und dass er diesen Zusammenhang
mit Hilfe der Theorie der Reihen \(p,q,r\) gefunden hat. Das
letztere wird zur Gewissheit, wenn man bedenkt, dass sich bei \textsc{Gauss} an keiner
Stelle auch nur der Versuch findet, für beliebige komplexe Werte der \(a,b\)
den Algorithmus des agM. \textls[100]{direkt} zu untersuchen; vielmehr geht er immer,\par
\clearpage\SourceLeaf{180}\ArithmeticHead{95}{R}{Vc. THEORIE DER MODULFUNKTION.}
\noindent wenn nicht die Voraussetzung positiver \(a,b\) gemacht wird, von dem Ansatz
aus:
\begingroup\makeatletter\tagsleft@true\let\veqno\leqno\makeatother
\begin{equation*}
a=hp(x)^2,\quad b=hq(x)^2,\quad c=(a^2-b^2)^{\frac12}=hr(x)^2,\tag{11}
\end{equation*}
\endgroup
\noindent wo man sich \(x\) als gegebenen Wert vom absoluten Betrage kleiner als Eins
zu denken hat (so z. B. auch in der Aufzeichnung 2) Werke X, 1, S. 210
und ebenso in 3) Werke X, 1, S. 218). Die Aufzeichnung 3) enthält (Werke
X, 1, S. 222—226) für die Reihen \(p,q,r\) eine vollständige Theorie der linearen
Transformation und den Zusammenhang dieser Theorie mit der Lehre von
den \textls[100]{quadratischen Formen}, also mit den \textls[100]{zahlentheoretischen} Untersuchungen
von \textsc{Gauss}. Er gibt daselbst für diese drei Reihen, die er als
Funktionen von \(t=-\frac{\log x}\pi\) in der Form schreibt
\begingroup\makeatletter\tagsleft@true\let\veqno\leqno\makeatother
\begin{equation*}
\left\{\begin{aligned}
p(t)&=1+2e^{-\pi t}+2e^{-4\pi t}+2e^{-9\pi t}+\cdots,\\
q(t)&=1-2e^{-\pi t}+2e^{-4\pi t}-2e^{-9\pi t}+\cdots,\\
r(t)&=2e^{-\frac14\pi t}+2e^{-\frac94\pi t}+2e^{-\frac{25}4\pi t}+\cdots,
\end{aligned}\right.\tag{12}
\end{equation*}
\endgroup
\noindent das Verhalten bei ganzzahliger linearer Transformation von \(t\),
\begingroup\makeatletter\tagsleft@true\let\veqno\leqno\makeatother
\begin{equation*}
t'=\frac{\alpha t-\beta i}{\delta+\gamma ti},\quad\alpha\delta-\beta\gamma=1,\tag{13}
\end{equation*}
\endgroup
\noindent in den sechs verschiedenen Fällen modulo 2, und es lässt sich auf Grund
dieser Formeln ein sehr plausibel erscheinender Gedankengang herstellen, der
sich in allen Einzelheiten den von \textsc{Gauss} aufgezeichneten Formeln anpasst
und eine durchaus einheitliche und in sich geschlossene Theorie des agM.
und der Modulfunktion liefert, die zum mindesten in ihren wesentlichen Zügen
der Ideenentwicklung von \textsc{Gauss} entsprechen dürfte.

Wir gehen wie \textsc{Gauss} in den Aufzeichnungen 2) und 3) von dem Ansatz
(11) aus, den wir aber in der Form
\begingroup\makeatletter\tagsleft@true\let\veqno\leqno\makeatother
\begin{equation*}
a=hp(t)^2,\quad b=hq(t)^2,\quad c=hr(t)^2\tag{11\,a}
\end{equation*}
\endgroup
\noindent schreiben wollen, wo also jetzt \(t\) eine gegebene Grösse mit positivem Realteil
bedeutet, und knüpfen nunmehr unmittelbar an die im Abschnitt IV b) (oben
S. 75ff.) skizzierten Entwicklungen an, nur schreiben wir auch die von zwei
Veränderlichen abhängenden \(\vartheta_{gh}\) statt mit den Argumenten \(\psi,x\) jetzt mit den
Argumenten \(\psi,t\)\footnote{Das schliesst sich an die in den neueren Darstellungen übliche Schreibweise (vergl. z. B. \textsc{Weber} a. a. O.\textsuperscript{102)}, S. 48) an, nur schreibt man mit \textsc{Jacobi} \(q\) (was dem Gaussschen \(x\) entspricht) als \(e^{\pi\tau i}\), also statt unseres \(t\), \(\tau=ti\); diese Abweichung können wir aber nicht vermeiden, da wir sonst nicht nur die Gaussschen Formeln, sondern auch die weiter unten wiederzugebenden Figuren abändern müssten. Es ist auch aus dem Grunde zweckmässiger, \(t\) statt \(x\) als die Variable zu betrachten, weil die Funktionen von \(t\) eindeutig abhängen, während die Abhängigkeit von \(x\) die mehrdeutige Potenz \(x^{\frac14}\) hat.}. Dann folgt aus den Gln. (29), (30) Abschnitt IV (S. 79)\par
\clearpage\SourceLeaf{181}\ArithmeticHead{96}{L}{L. SCHLESINGER, ÜBER GAUSS’ FUNKTIONENTHEORETISCHE ARBEITEN.}
\begingroup\makeatletter\tagsleft@true\let\veqno\leqno\makeatother
\begin{equation*}
\left\{\begin{gathered}
a_n=hp(2^nt)^2,\quad b_n=hq(2^nt)^2,\quad c_n=hr(2^nt)^2,\quad(n=1,2,\ldots)\\
h=M(a,b),
\end{gathered}\right.\tag{14}
\end{equation*}
\endgroup
\noindent wodurch ein eindeutig bestimmter Wert des agM. zwischen \(a\) und \(b\) fixiert
ist. — Nun finden wir im Art. 6 der Aufzeichnung 3) die Formel für die
lineare Transformation der Funktion \(\vartheta_{00}(\psi|t)\), deren Beweis \textsc{Gauss} im Art. 1
der Notiz aus dem Handbuch 16, Bb (Werke X, 1, S. 287ff., aus dem Jahre
1808) gegeben hat:
\begingroup\makeatletter\tagsleft@true\let\veqno\leqno\makeatother
\begin{equation*}
\vartheta_{00}\left(\frac\psi{ti}\middle|\frac1t\right)=e^{\frac{\pi\psi^2}t}t^{\frac12}\vartheta_{00}(\psi|t),\tag{15}
\end{equation*}
\endgroup
\noindent wo\footnote{Vergl. z. B. \textsc{Weber} a. a. O.\textsuperscript{102)}, S. 77, Gln. (11).}\ \(t^{\frac12}\) so zu wählen ist, dass sein reeller Teil positiv ausfällt. Aus dieser
Formel und den sofort daraus herzuleitenden für \(\vartheta_{01},\vartheta_{10}\) ergeben sich für
\(\psi=0\) die Formeln
\begingroup\makeatletter\tagsleft@true\let\veqno\leqno\makeatother
\begin{equation*}
p\left(\frac1t\right)=t^{\frac12}p(t),\quad q\left(\frac1t\right)=t^{\frac12}r(t),\quad r\left(\frac1t\right)=t^{\frac12}q(t),\tag{16}
\end{equation*}
\endgroup
\noindent die \textsc{Gauss} mit den aus den Reihenentwicklungen (12) unmittelbar folgenden
\begingroup\makeatletter\tagsleft@true\let\veqno\leqno\makeatother
\begin{equation*}
p(t+i)=q(t),\quad q(t+i)=p(t),\quad r(t+i)=e^{\frac{\pi i}4}r(t)\tag{17}
\end{equation*}
\endgroup
\noindent im Art. 6 von 3) (Werke X, 1, S. 223) aufgezeichnet hat\begingroup\renewcommand{\thefootnote}{131a)}\footnote{\textsc{Gauss} schreibt a. a. O. diese Formeln direkt für \(p(t)\) hin, für die beiden andern Reihen nur in der Form, dass er sich \(q,r\) durch die Gleichungen
\[
q(t)=2p(4t)-p(t),\quad r(t)=\tfrac12p(\tfrac14t)-\tfrac12p(\tfrac14t+i)
\]
definiert denkt. Übrigens finden sich dieselben Formeln auch in der weiter unten wiederzugebenden Aufzeichnung aus der Scheda An, die schon aus dem Jahre 1805 stammt.}\endgroup\addtocounter{footnote}{-1}. Aus den Gleichungen
(16) folgt mit Rücksicht auf \(M\bigl(p(\tfrac1t)^2,q(\tfrac1t)^2\bigr)=1\) (siehe oben Gl.
(30), Abschn. IV, S. 79) die Formel
\begingroup\makeatletter\tagsleft@true\let\veqno\leqno\makeatother
\begin{equation*}
\frac1t=M(p(t)^2,r(t)^2)=\frac{M(a,c)}{M(a,b)},\tag{18}
\end{equation*}
\endgroup
\noindent wodurch auch eine ganz bestimmte Determination von \(M(a,c)\) festgelegt ist,
von der wir sagen, dass sie mit dem oben bestimmten \(M(a,b)\) \textls[100]{zusammen-}\par
\clearpage\SourceLeaf{182}\ArithmeticHead{97}{R}{Vc. THEORIE DER MODULFUNKTION.}
\HistoricalFoot{X 2 Abh. 2.}{13}
\noindent\textls[100]{gehört}. — Die Folgerung (18) aus den Gln. (16) bezw. aus der Gl. (15) hat
\textsc{Gauss} in der oben genannten Notiz aus Handbuch 16, Bb, gezogen und mit
den Worten (Werke X, 1, S. 289) »das agM. zwischen \(p(x)^2\) und \(r(x)^2\) ist \(\dfrac\pi{\log\dfrac1x}\)«
aufgezeichnet. Sie bildet eine wesentliche Ergänzung der auf die Reihen
\(p,q,r\) bezüglichen Formeln (28), (29), (30) Abschn. IV (oben S. 78, 79) und
liefert, wie wir jetzt zeigen wollen, eine ebensolche auch für die Formeln
(31) bis (36) desselben Abschnitts (oben S. 80ff.). Wenn nämlich die durch
Gl. (31) Abschn. IV definierte Funktion
\[
\operatorname{Sl}(\varphi)=\frac{q(t)}{r(t)}\frac{\vartheta_{01}(\psi|t)}{\vartheta_{00}(\psi|t)}
\]
\noindent die Differentialgleichung (vergl. (35) a. a. O.)
\[
\frac{dS}{d\varphi}=\sqrt{(1-S^2)(1+\mu^2S^2)},\quad\mu^2=\frac{r(t)^4}{q(t)^4},
\]
\noindent befriedigt, so können wir für \(\mu=\operatorname{tg}v\), wegen \(p^4-q^4=r^4\), setzen:
\begingroup\makeatletter\tagsleft@true\let\veqno\leqno\makeatother
\begin{equation*}
k=\sin v=hr(t)^2,\quad k'=\cos v=hq(t)^2,\quad1=hp(t)^2,\tag{19}
\end{equation*}
\endgroup
\noindent woraus sich unmittelbar \(h=M(1,\cos v)\) ergibt. — Die Gleichung (18) liefert
dann
\begingroup\makeatletter\tagsleft@true\let\veqno\leqno\makeatother
\begin{equation*}
t=\frac{M(1,\cos v)}{M(1,\sin v)},\tag{20}
\end{equation*}
\endgroup
\noindent und wenn wir nunmehr \(\varpi'=t\varpi\) setzen, so wird
\stepcounter{footnote}
\begingroup\makeatletter\tagsleft@true\let\veqno\leqno\makeatother
\begin{equation*}
\varpi=\pi q(t)^2=\frac{\pi\cos v}{M(1,\cos v)},\quad
\varpi'=\pi tq(t)^2=\frac{\pi\sin v}{M(1,\sin v)}\text{\textsuperscript{132)}}.\tag{21}
\end{equation*}
\endgroup
\footnotetext{Ziehen wir noch die Gleichung (7) heran so finden wir
\[
\begin{gathered}
\int_0^\pi\frac{du}{\sqrt{1-\sin^2v\cos^2u}}=\pi p(t)^2,\\
\int_0^\pi\frac{du}{\sqrt{1-\cos^2v\cos^2u}}=\pi tp(t)^2,
\end{gathered}
\]
welche Gleichungen \textsc{Gauss} neben den aus (19) folgenden
\[
q(t)=p(t)\sqrt{\cos v},\quad r(t)=p(t)\sqrt{\sin v}
\]
im Art. 9 der Notiz 3) (Werke X, 1, S. 226, unten, rechts vom Strich) aufgezeichnet hat.}
\noindent Die hier auftretenden Funktionen von \(t\)
\[
\mu^2=\frac{r^4}{q^4},\quad k^2=\frac{r^4}{p^4},\quad k'^2=\frac{q^4}{p^4}
\]
\noindent stehen weiterhin im Mittelpunkt von \textsc{Gauss'} Untersuchungen, insbesondere\par
\clearpage\SourceLeaf{183}\ArithmeticHead{98}{L}{L. SCHLESINGER, ÜBER GAUSS’ FUNKTIONENTHEORETISCHE ARBEITEN.}
\noindent die letzte \(k'^2\), die wir als Modulfunktion im engeren Sinne zu bezeichnen gewohnt sind. —
